% \section{Combined Conceptual Deductive Analysis}
\section{Deductive Design of Conceptual Architecture}

Based on the deep knowledge gained by the previous sections,
it is now possible to create a combined conceptual approach to address the issues
described in the \nameref{Introduction} and the \nameref{ProblemStatement} sections.
This approach will be used as a guideline for the prototype implementation int the chapter \nameref{ArtifactDevelopment}.

In contrast to the sections above which went from the basic principles towards the higher levels of abstraction,
this chapter will start at the high level goals, their implications and what will be needed on the lower levels
in order to accomplish these goals.

\subsection{Goals}
\label{sec:goals}

Based on the \nameref{ProblemStatement}, the goals of the system can be summarized as follows:

% Original problem statement:
% While low latency high performance computing has been a focus of the HPC as well as the
% CC community, resulting in systems such as JupyterHub or Dask, these systems are inherently
% ephemeral and do not provide a way to reproduce results on a HPC scale. Therefore we need to
% solve the following problems:
% • P1: The GitOps first approach with introduces a latency between the changes made and
% the deployment of the new workflow. This latency prevents a real-time exploitive workflow
% and limits the usability for the interactive exploration of the data.
% • P2: Existing solutions for interactive workloads are inherently ephemeral and do not pro-
% vide a way to reproduce the results on a HPC scale


The designed system should provide the user with an easy-to-use interface,
capable of interactively exploring \ac{HPC} scale data.
The system should be able to run on cloud infrastructure, be able
to allocate the necessary resources, as well as scheduling and executing
the computation within a reasonable time frame.

% Yeah honestly it was not realistic to expect that I would be able to build a system which 
% could run HPC on the cloud just like that, while also seamlessly autoscaling and being cheap lol
% considering that this needed to be done within two months while writing this thesis at the same time.
% I got so close tho :(

% Commented out as this is not just unrealistic, but misses the original problem statement:
% On the highest level, we want a system which provides the user with a 
% computational notebook, which is able to execute \ac{HPC} scale computations seamlessly on cloud infrastructure.
% The system should be able to automatically provide the necessary resources on demand
% and in a reasonably cost effective way.

% Notes: its a given that it is supposed to run on cc infr??
% what would the advantages even be about it running on cloud infrastructure in the first place?

\subsection{Implications}

In order to achieve the goals, the system must be able to manage the following aspects, in order of abstraction:

\begin{itemize}
    \item \textbf{Programming Language Support:} in order to enable \ac{HPC} scale computation for the user,
          the programming language must support parallel computation natively, or at least provide a library which does.
          As programming languages such as Python are not natively able to support parallel computation paradigms
          such as described in the \nameref{sec:programming_models} section.
          The language must not just support parallel computation, but also be user-friendly in order to let the user focus on
          exploring the data and not spend their time on the intricacies of parallel computation.
    \item \textbf{Notebook Environment:} The system should support a Notebook environment, allowing users to
          facilitate the communication between the client and the compute cluster, requiring an endpoint being able to submit
          the code to the cluster and receive the results.
    \item \textbf{Compute Cluster:} The system must be able to schedule and execute the computation within a reasonable time frame.
          Requiring a system capable of managing the resources and scheduling the computation.
    \item \textbf{Platform Agnostic:} As the system is supposed to run on cloud infrastructure,
          whether it is on-premise or in the cloud, the system must be able to run cloud natively.
\end{itemize}


\subsection{Suggested Architecture}
\label{suggested_architecture}

Based on the requirements outlined above a hypothetical software can be suggested
to fulfill the requirements. While the implementation of this in the following chapter will additionally
constrain by the available resources and time frame here we can outline a system which
is not limited by these circumstances.

\subsubsection{Service Layer Architecture}

This section will cover the highest level of abstraction,
the service layer architecture, which will be the most important part of the system for the user.

\textbf{Client}
\label{asp:client}

With the usage of a Compute Notebook being an essential part of exploratory data analysis,
and as discussed in the \nameref{sec:EC} section, Jupyter Notebook is the most popular choice
and will bring a familiar interface to scientists and domain experts, who had little experience with \ac{HPC} before.

In  this case we will make use of the Jupyter Server Client architecture,
having the user only run the client locally while the server will be running on the cloud infrastructure.

\textbf{Language}
\label{asp:language}

The language of choice is \textit{Arkouda}, a highly parallel framework for python,
specifically build for enabling data scientists to perform parallel computation on terrascale \ac{HPC} systems on a high level of abstraction
and interactive speeds\footcite{enthoughtArkoudaTerascaleData2020}.

It is build on top of the \textit{Chapel} programming language, which itself is a parallel programming language
specifically designed on \ac{HPC} systems, and instead of being a set of libraries on top of a general purpose language
it is build from the ground up to work on \ac{PGAS} systems, which is a perfect fit for the use case of this system as
will be discussed later on in the \nameref{sec:suggarch_network} section.

Another option could be \ac{MPI}, as it is almost as user-friendly as Python,
but as Python is a more popular language and is more user-friendly, it would be a better fit for the system.

% - Python a quite popular programming language not just with data scientist but also domain experts.
% - Needs to support parallel computation, as the language itself does not support it natively
%   - therefore the arkouda ?dialect? is used as it is especially greared towards this kind of computation.
%   - uses a chapel interpreter to execute the code and to make use of distibuted compute hardware

\textbf{Cluster}
\label{asp:cluster}

While the development of containers was primarily driven by \ac{CC} needs and those two subjects
how become tightly linked, \ac{CC} is much more than just containers.
Complete \ac{CC} systems such as \textit{OpenStack} don't limit themselves to just containers,
but also provide virtual machines as well as bare metal provisioning to provide compute resources\footcite{jacksonOpenStackCloudComputing2013}.

For the reasons discussed in the \nameref{sec:Cloud_Computing} section, virtualization is not a good fit,
as in the best case we would still emulate a full operating system, which would be a huge overhead.

Bare metal, which is the practice of simply foregoing any kind of virtualization or containerization
and just running the code on the hardware itself, at the cost of not being able to subdivide the resources of a node\footcite{canonicalMetalService}.
This is done by creating a custom operating system image which is installed via a network boot.
This means that provisioning a new node is extremely slow in comparison to starting an isolated container,
but also ensures 100\% of the resources are available to the user and no overhead is introduced by the virtualization layer.
The specific use case this could fit is if the entire nodes compute resources would be taken up by a single service
anyhow but most importantly the resource demand is very stable and predictable.

While the former is true and would make a compelling case for the compute-nodes,
but especially in the context of exploratory computation the resource demands are highly dependent on the user
and can change rapidly, which would make the latter a bad fit for the system.

The best fit for the use case is a container based system, as spawning containers is as fast as starting a new process,
while bringing isolation, resource subdivision and a rich ecosystem of tools and libraries to the table.

This cluster will be hosting the entire system, from the compute services doing our \ac{HPC} computation,
to the exposed services such as the Jupyter Server, which will be able to communicate with the client.

% - Kubernetes as a container orchestration system vs swarm
% - talk about how KVM could also be used, but with the drawback of beings slower to provision and have a higher overhead
% - also a hardware provisioning system could be used such as OpenStack but this is not just much slower
% but also prevents a secure way of sharing the resources between mutlitple users and is a much less flexible system


\textbf{Jupyter Server}
\label{asp:jupyter_server}

The Jupyter Server web application will be running on the cluster being the developers access point to the system.
It will be able to communicate with the client, submit the code to the compute cluster and request the necessary resources.

While there is also a multi-user Jupyter Server called JupyterHub, which allows multiple users to share the same server, we will still be using the single user instance\footcite{jupyterhub-docsJupyterHub}.
If more than one user wants to use the system at a given time, instead of sharing the same environment they each will get their very own
complete environment, which will be spawned on demand and destroyed after the user is done with it.
This will ensure not just a clean environment but also prevent any kind of interference between the users.

At least one Jupyter Server will always be running on the cluster, ready to accept new connections from the client.


\textbf{Compute Services}
\label{asp:compute_services}

Now that the user interface and its functionality have been discussed, the next step is to examine the compute services.
These services are responsible for executing the code submitted by the user on the compute nodes.
As we have already discussed, the language of choice is Arkouda.
It is specifically designed for enabling interactive programming and parallel computation.
Arkouda splits the computation from the developer's machine to the compute nodes,
allowing for efficient and scalable processing.
Its infrastructure is split into three parts\footcite{arkouda-docsStartupArkoudaDocumentation}:


\begin{itemize}
    \item \textbf{Client:} This can be any python interpreter running on the users machine, which will be able to communicate with the server.
    \item \textbf{Server:} This is the bootstrap server, it is the entry point for the client and bootstraps the multi-locale cluster.
    \item \textbf{Locales:} These are the workers, at least one of then will be running on each node that contributes to the computation.
\end{itemize}

For the system to be build, both the server and an arbitrary amount of locales will be running on the cluster,
with the server establishing the initial connection to the locales through \ac{SSH}, setting up
the communication in order to set up the \ac{PGAS} space.

The server will be on standby until a connection request is made by a client and will then submit the
processing request to the locales through the interconnect.

\textbf{Resource Provisioning / Scaling}
\label{asp:resource_provisioning}

While in traditional \ac{HPC} systems the resources are statically allocated and the user has to request the resources beforehand in order to be scheduled,
the system we are building is supposed to be more dynamic and able to scale the resources on demand.
Similarly to the time google initially developed the Borg system in order to handle both their services as
well as their batch jobs, as described in the \nameref{sec:orchestration} section, it is to be expected
that this system  might simultaneously handle other services as well as interactive computation,
in order to make the best use of the hardware invest.

While an ideal solution would be able to scale the available resources for each single submission of a code cell,
the way the Arkouda system is designed, the size of available nodes needs to be known beforehand.
This is due to the fact that it relies on the \ac{PGAS} model (see \nameref{sec:gasnet}) which requires
the memory to be shared between the nodes, if we now started to add or remove nodes during the computation
the memory space would need to be resized each time which would come with a large overhead.
A possible solution to this problem could be the use of a system such as \ac{FAM} which is able to provide an external memory space
which is able to survive program termination and is made as a shared memory space between the nodes\footcite{keetonOpenFAMAPIProgramming2019a}.

An alternative approach could be to work on the scale of a session instead of a single code cell,
where the user would be provided with a fixed amount of resources for the entire session,
depending on the utilization of the system, the user could be provided with more or less resources.
This would be a more stable approach, as the user would not be interrupted by the system scaling up or down,
but would also not be able to make use of the full potential of the system.

Depending on circumstances, another factor in this case could be Kubernetes ability to
over provision resources.
Over provisioning in this case means that while one node is technically bound up in a session,
it could still be used by a different session, while the cores are sitting idle.
In classical \ac{HPC} systems this would not be possible as the compute request
has been scheduled long before, but as exploitive computing is inherently burst,
sitting idle while the user is thinking or writing code, these nodes could be used by other users.
But this system is inherently a trade-off between the user experience and the utilization of the system and
depends entirely on the specific use case.

\subsubsection{Cluster Level Architecture}

This section will cover the lower level of the system architecture,
the sections which are not directly visible to the user but are essential for the system to work.

\textbf{Operating System}
\label{asp:operating_system}

Using containerization as the main method of resources isolation makes the
choice of Linux as a base an imperative.
As discussed in the section \nameref{fig:containerization} containers are
build around features specifically created in the Linux kernel.
While other operating systems do technically support some form of containerization,
either through creating a Linux virtual machine or creating
a custom containerization system, these are not as mature as the Linux based solutions
\footcite{mcgrathServerlessComputingDesign2017}.
Aside from specific feature support regarding containers,
Linux has become the de facto standard for \ac{HPC} systems (\ref{sec:operating_systems})
as well as cloud infrastructure \footcite{stevenvaughan-nicholsCanInternetExist},
due to its open source nature and the ability to customize it to the specific needs of the industry.

While the choice of operating system is clear, the choice of which \gls{distribution} is much less so.
Essentially there are four kinds of distributions which could be used,
while they will have little functional difference for the system build on top of them,
as containerization abstracts the underlying system away, they will have a large impact on the maintainability of the system,
the availability of support and the ease of development of the system itself.

\begin{itemize}
    \item \textbf{Cloud Managed Distributions:} The large hyperscalers such as AWS\footcite{awsOpenSourceKubernetesVerteilungAmazonEKSVerteilungAmazon}, Google Cloud\footcite{googleNodeImagesGoogle}
          or Azure\footcite{burtMicrosoftAzureLinux} all provide their own distributions within their cloud environments, which are specifically tailored to their needs.
          These are specifically geared towards the cloud environment and are not meant to be used outside it.
    \item \textbf{Enterprise Distributions:} These are distributions which are specifically tailored to the needs of the enterprise customers,
          running their own data centers. They provide a fully featured out-of-the-box solution, with professional support and a large ecosystem of tools.
          Examples of this are \textit{Red Hat OpenShift}\footcite{redhatRedHatOpenShift} or \textit{VM Ware Tanzu Kubernetes Grid} \footcite{vmwareVMwareTanzuKubernetes2020}.
          These systems are costly but provide a high level of support and are specifically tailored towards a specific way of administration.
    \item \textbf{Specific \ac{FOSS} Distributions:} These are distributions which are build around the idea of being open source and free to use
          and are extendable by the user. Examples of this are \textit{Suses Rancher}\footcite{rancherInnovateEverywhere} or the community edition of \textit{OpenShift} called \textit{okd} \footcite{innesOKDIo}.
          While some offer a professional support option, they can be used and extended by the user as they see fit.
    \item \textbf{Generic \ac{FOSS} Distributions:} These are distributions which are not specifically tailored towards any specific use case, but as the features of
          kubernetes are build around the kernel which is shared by all distributions, they can be used as well.
          But will require a manual installation of the necessary software and will not have any kind of professional support. This includes distributions such as \textit{Debian}\footcite{debianDebianUniversalOperating} or \textit{Ubuntu}\footcite{canonicalGetUbuntuServer}.
\end{itemize}

As the system we will be building is going to be highly specific and will be running software which is out of the ordinary,
from a cloud infrastructure perspective, the utility of outside support might be limited and the flexibility of the system might be more important.
Therefore, the choice between a generic \ac{FOSS} distribution or a specific \ac{FOSS} distribution depends on factors such as the need for ready-made integrations and GUIs,
as well as the administrators' familiarity with different systems.
In this case, a Linux-based operating system would be the most suitable choice.

% as containers as described are based around the Linux kernel specifically.

% Disussion of dedicated Containeroriented os vs gernric linux distro
% but anywhow one with professional support.
% Opensift? Rancher? 

% Or debian with manual installation.

\textbf{Container Runtime}
\label{asp:container_runtime}

The choice of container runtime is a crucial one, as the container runtime is the software doing the "containerizatoin" --
the isolation of the resources, the bootstrapping of the container and the execution of the code as described in the \nameref{sec:containerization} section.
The \ac{CNCF} recognizes 19 different container runtimes \footcite{cncf[cloudnativecomputingfoundation]CNCFLandscape},
the most relevant in our case are \textit{containerd}, \textit{\ac{CRI-O}} and \textit{singularity}.
\textit{containerd} is the runtime used by docker, which has a high market share and is well known,
due to it being the one of the first container runtime to be widely adopted both of the others are improvements or specializations of it.

\textit{\ac{CRI-O}} is a lightweight container runtime specifically designed for Kubernetes and is build around the \ac{OCI} standard,
being more streamlined towards large scale \ac{CC} systems\footcite{crioCrioCrio2024} but for the application of this system would
work mostly interchangeably with \textit{containerd}.

\textit{Singularity} on the other hand is a container runtime specifically designed for \ac{HPC} systems\footcite{singularity-docsIntroductionSingularitySingularity},
which are trying to run as close to the metal as possible, and to support the specific needs of \ac{HPC} systems
such as direct access to external accelerators such as \ac{GPU}s and \ac{FPGA}s, support for \ac{HPC} specific interconnects\footcite{uni.luULHPCTutorial} (see \ref{sec:networking_interconnects}) and the ability to run
in \gls{rootless} mode\footcite{singularity-docsFakerootFeatureSingularity}.


\textbf{Container Image}
\label{sec:container_image}

The container image will be an extension of the base image which
can be fund in the  Arkouda repo \footcite{arkouda-repoArkoudacontribArkoudadockerMain},
which itself is based upon the chapel container image \footcite{chapelChapelProfileDocker}.
It makes chapel and Arkouda available to be used in a containerized environment.


% \textbf{Scheduling System}
\textbf{Network Abstraction}
\label{asp:network_abstraction}

The network abstraction layer creates a homogenous network between the different containers even if they are running on different nodes
This makes the code running in the containers much easier to scale, manage and write.
This can become quite complex involving multiple layers of abstraction and
nested protocols such as described in the \nameref{sec:K8sNetworking} section.

But while these systems make the network much easier to manage,
they are a trade-off between usability and performance, as the layers of abstraction
such as \ac{VXLAN} encapsulation or \ac{IP} in \ac{IP} tunneling come with a performance overhead.
While this overhead is negligible for most use cases,
in \ac{HPC} systems where the network is the bottleneck it can be a significant factor.

While the most common network overlays in Kubernetes such as \textit{Flannel} and \textit{Calico},
support both \ac{VXLAN} and \ac{IP} in \ac{IP} tunneling\footcite{tigeraOverlayNetworkingCalico}, as this is needed to communicated
across subnets in a cloud environment, but for performance reasons
they also support foregoing a virtual network and using the physical network directly\footcite{flannel-repoFlannelDocumentationBackends}

But as Arkouda and many other \ac{HPC} systems are build around the \ac{RDMA} as an available feature,
we also need an available network which supports \ac{RDMA} as well. This can be
done by creating an additional pod level network dedicated the \ac{RDMA} communication between the nodes,
using \textit{Multus}\footcite{multusGitHubK8snetworkplumbingwgMultuscni} and then the specific network plugin
for the \ac{RDMA} capable network card\footcite{nvidiaNetworkOperator}.

\pagebreak

\subsubsection{Hardware Level Architecture}

\textbf{Compute Nodes Hardware}
\label{asp:compute_node_hardware}

Regarding the hardware of the nodes making up the cluster,
many rules which are true for \ac{HPC} systems described  in the \nameref{sec:hpc_infrastructure} section are also true for this system.
For this system, it is important to consider factors such as the need for low latency and high bandwidth in order to facilitate efficient communication in \ac{HPC} systems.
Additionally, the computational power should be as high as reasonably possible to handle large-scale computations effectively.

In order to handle smaller scale computations which fit into a single node,
a High core count CPU would be a good fit as well as a large amount of memory,
which will also come in handy for the larger scale computations, as part of the \ac{PGAS}.

One advantage of Kubernetes is the ability to use more heterogeneous hardware than in traditional \ac{HPC} systems.
This allows for the use of nodes of different sizes and capabilities, optimizing the system for the specific use case while
still using the same software stack and administration tools.

Overall this part of the system is mostly governed by the budget available,
as more powerful hardware will be able to handle larger computations faster.
And if the use case requires specific hardware, such as \ac{GPU}s or \ac{FPGA}s,
the system should be able to accommodate these as well.

\textbf{Network Interconnect}
\label{sec:suggarch_network}

For the underlying interconnect, it is crucial to use a network that supports \ac{RDMA}. However, in cloud environments, the most common network interconnects are simple Ethernet networks,
which do not have built-in support for \ac{RDMA}.
While Arkouda is able to use \ac{PGAS} which can use any kind of network,
by wrapping the communication in a higher level protocol, this is primarily for proof of concept and small scale computations,
as resolving the communication needs to be done in the software stack,
instead of being able to rely on the network hardware to do it for you.

For this reason it would be best to use a network which supports \ac{RDMA} natively,
such as InfiniBand or SlingShot, which are able to provide the necessary low latency and high bandwidth
needed for \ac{HPC} systems\footcite{desensiInDepthAnalysisSlingshot2020a}.

Topologically the network would differ from traditional cloud network topologies,
which are usually less interconnect than \ac{HPC} systems,
as a lot of the communication comes from the outside of the cluster and not between the nodes themselves.
Our network would instead be ideally a full bisectional bandwidth network,
where every node is connected to every other node, in order to provide the lowest latency possible.
But more realistically a fat tree network topology would be used, just as discussed in the \nameref{sec:network_topology} section.

% frontend: plain jupyter notebook
% communication: jupyter protocol
% backend:
%     - infrastructure
%       - hardware
%         - 20 node cluster
%         - storage system
%       - interconnect
%         - ethernet without RDMA
%           - evaluation of possobilies
%     - software
%       - container orchestration system: Kubernetes
%       - workload management system: Kubernetes scheduler
%       - compute software : Arkouda/Chapl
%       - storage software: Minio
%       - neworksoftware: 


% % TODO:
% Talk about how Arkouda is an easy to use library able to provide the user with the necessary tools to perform parallel computation.
% Talk about how Arkouda supprots PGAS under any kind of network infrastructure, even without RDMA capale interconnect
% But if RDMA is there how container could still be used to provide the user with the necessary tools to perform parallel computation
% (talking about namespace isolaton ...)
% Talk about how the containers can make use of the underlying CPU architecture even in a cloud system.
% Talk about how the container are able subdevide the resources of the underlying hardware even for multiple users.
% Talk about alternative Kubernetes scheduling systems interlacing deployments and jobs.